\subsection{Hoja de Ruta (\textit{Roadmap}): Cronograma, Entregables y Recursos}
\label{subsec:roadmap_detallada}

La hoja de ruta consolidada de la Tabla~\ref{tab:roadmap_cci} se desagrega en esta subsección en un plan ejecutable.
Cada fase se descompone en paquetes de trabajo con una duración estimada, se asocia con hitos verificables y con los cuatro
entregables principales de la arquitectura, y se completa con la asignación de responsables y esfuerzo. El plan se expresa en
meses relativos al inicio del programa (M1), de modo que pueda calendarizarse cuando se defina la fecha de arranque, y mantiene
los horizontes de 6, 12 y 24 meses: la Fase~1 se ejecuta entre M1 y M6, la Fase~2 entre M7 y M12 y la Fase~3 entre M13 y M24.

En total, el plan contiene 21 paquetes de trabajo (siete por fase), 14 hitos, cuatro entregables principales y 14 productos de
trabajo asociados. El detalle completo, junto con las fórmulas de cálculo de fechas y esfuerzos, se encuentra en el libro
complementario \texttt{Hoja\_de\_Ruta\_CCI.xlsx}.


\subsubsection{Cronograma de hitos y estimaciones de tiempo}

La Tabla~\ref{tab:cronograma_wp} presenta los paquetes de trabajo de cada fase con su mes de inicio, mes de finalización,
duración estimada y dependencias, y la Figura~\ref{fig:gantt_cci} los muestra sobre una línea de tiempo común. Las duraciones
son estimaciones de planificación propuestas para esta hoja de ruta y deben confirmarse con la Dirección TIC y los equipos
involucrados antes de comprometer fechas.

\begingroup
\scriptsize
\setlength{\tabcolsep}{3pt}
\renewcommand{\arraystretch}{1.15}
\begin{longtable}{
        >{\raggedright\arraybackslash}p{0.08\textwidth}
        >{\raggedright\arraybackslash}p{0.43\textwidth}
        >{\centering\arraybackslash}p{0.07\textwidth}
        >{\centering\arraybackslash}p{0.07\textwidth}
        >{\centering\arraybackslash}p{0.07\textwidth}
        >{\raggedright\arraybackslash}p{0.18\textwidth}}
    \caption{Cronograma de paquetes de trabajo de la hoja de ruta del CCI}
    \label{tab:cronograma_wp}                                                                                                                              \\
    \toprule
    \textbf{ID} & \textbf{Paquete de trabajo}                                   & \textbf{Inicio} & \textbf{Fin} & \textbf{Dur.} & \textbf{Depende de}     \\
    \midrule
    \endfirsthead
    \toprule
    \textbf{ID} & \textbf{Paquete de trabajo}                                   & \textbf{Inicio} & \textbf{Fin} & \textbf{Dur.} & \textbf{Depende de}     \\
    \midrule
    \endhead
    \midrule
    \multicolumn{6}{r}{\footnotesize Continúa en la página siguiente}                                                                                      \\
    \endfoot
    \bottomrule
    \endlastfoot
    \multicolumn{6}{l}{\textbf{Fase~1 · Gobierno de identidades y accesos (M1--M6)}}                                                                       \\
    \midrule
    WP1.1       & Inventario de identidades y diseño del modelo RBAC            & M1              & M3           & 3             & —                       \\
    WP1.2       & Activación progresiva de MFA en las consolas del CCI          & M3              & M6           & 4             & 1.1                     \\
    WP1.3       & Consolidación de RBAC y segregación de funciones              & M4              & M6           & 3             & 1.1                     \\
    WP1.4       & Registro automático de decisiones                             & M2              & M5           & 4             & —                       \\
    WP1.5       & Definición de fuentes autoritativas de datos de flota         & M1              & M4           & 4             & —                       \\
    WP1.6       & Segmentación y protección perimetral inicial                  & M4              & M6           & 3             & —                       \\
    WP1.7       & Verificación de entregables y cierre de la Fase 1             & M6              & M6           & 1             & 1.2, 1.3, 1.4, 1.5, 1.6 \\
    \addlinespace[3pt]
    \multicolumn{6}{l}{\textbf{Fase~2 · Telemetría confiable y trazabilidad (M7--M12)}}                                                                    \\
    \midrule
    WP2.1       & Identidad para dispositivos ITS                               & M7              & M9           & 3             & 1.7                     \\
    WP2.2       & Autenticación de mensajes GPS/AVL en la pasarela de ingreso   & M9              & M11          & 3             & 2.1                     \\
    WP2.3       & Validación de plausibilidad y cuarentena de tramas inválidas  & M9              & M11          & 3             & 2.1                     \\
    WP2.4       & Integración segura del bus de eventos y fuentes geoespaciales & M7              & M11          & 5             & 1.5                     \\
    WP2.5       & WORM, sellado de tiempo y encadenamiento de hashes            & M7              & M11          & 5             & 1.4                     \\
    WP2.6       & Correlación en SIEM y monitoreo centralizado                  & M9              & M12          & 4             & 2.2, 2.5                \\
    WP2.7       & Verificación de entregables y cierre de la Fase 2             & M12             & M12          & 1             & 2.1 a 2.6               \\
    \addlinespace[3pt]
    \multicolumn{6}{l}{\textbf{Fase~3 · Continuidad, redundancia híbrida y recuperación (M13--M24)}}                                                       \\
    \midrule
    WP3.1       & Diseño de la arquitectura híbrida redundante                  & M13             & M15          & 3             & 2.7                     \\
    WP3.2       & Replicación de servicios y datos CCI–GCP                      & M16             & M20          & 5             & 3.1                     \\
    WP3.3       & Balanceo y conmutación automática (failover)                  & M18             & M21          & 4             & 3.1, 3.2                \\
    WP3.4       & Consolidación del sitio alterno                               & M16             & M21          & 6             & 3.1                     \\
    WP3.5       & Formulación y aprobación de DRP y BCP                         & M13             & M18          & 6             & 2.7                     \\
    WP3.6       & Pruebas de recuperación y failover                            & M21             & M23          & 3             & 3.3, 3.4, 3.5           \\
    WP3.7       & Verificación de entregables y cierre de la Fase 3 (TO-BE)     & M24             & M24          & 1             & 3.2 a 3.6               \\
\end{longtable}
\endgroup

{\footnotesize\textit{Nota.} Los meses son relativos al inicio del programa (M1); Dur. es la duración en meses. En la columna «Depende de» se indican los
    paquetes de trabajo (sin el prefijo WP) cuyos resultados se requieren como insumo; algunas actividades se solapan
    parcialmente con su predecesora para permitir entregas incrementales. Elaboración propia a partir de las vistas de
    Project e Implementation and Migration de ArchiMate \parencite{archimate2022,modeloArchimateCCI2026} y de los horizontes de 6, 12 y 24 meses \parencite{fase1cci2026}.\par}


La Fase~1 (M1--M6) concentra los controles de menor dependencia de transformaciones de infraestructura. El inventario de
identidades y el diseño del modelo RBAC (WP1.1) abren la fase y habilitan tanto la activación progresiva de MFA (WP1.2) como la
consolidación de RBAC y de la segregación de funciones (WP1.3). En paralelo se automatiza el registro de decisiones (WP1.4), se
definen las fuentes autoritativas de datos de flota (WP1.5) y se inicia la segmentación y protección perimetral (WP1.6). La secuencia
que condiciona el cierre de la fase es WP1.1, WP1.2 y la verificación de WP1.7.

La Fase~2 (M7--M12) protege la telemetría y la trazabilidad. La identidad de los dispositivos ITS (WP2.1) es el primer paso,
porque la autenticación de mensajes GPS/AVL (WP2.2) y la validación de plausibilidad con cuarentena de tramas inválidas (WP2.3)
dependen de ella. La integración segura del bus de eventos (WP2.4) y el fortalecimiento de la bitácora con WORM, sellado de tiempo y
encadenamiento de hashes (WP2.5) avanzan en paralelo desde M7, y la correlación en SIEM (WP2.6) se apoya en los resultados de
ambos frentes. La secuencia que condiciona el cierre es WP2.1, WP2.2, WP2.6 y WP2.7.

La Fase~3 (M13--M24) consolida la continuidad operativa. El diseño de la arquitectura híbrida redundante (WP3.1) precede a la
replicación de servicios y datos (WP3.2), al failover automático (WP3.3) y a la consolidación del sitio alterno (WP3.4), mientras que
el DRP y el BCP se formulan en paralelo (WP3.5). Las pruebas de recuperación y failover (WP3.6) solo pueden iniciar cuando
conmutación, sitio alterno y planes están disponibles, por lo que la secuencia WP3.1, WP3.2, WP3.3, WP3.6 y WP3.7 determina la fecha de
llegada al estado TO-BE. Las pruebas son periódicas, por lo que continúan en operación después de M24.

\begin{landscape}
    \begingroup
    \definecolor{ganttF1}{HTML}{2E75B6}
    \definecolor{ganttF2}{HTML}{548235}
    \definecolor{ganttF3}{HTML}{C55A11}
    \definecolor{ganttHito}{HTML}{C00000}
    \begin{figure}[H]
        \centering
        \caption{Diagrama de Gantt de la hoja de ruta del CCI (meses relativos al inicio del programa)}
        \label{fig:gantt_cci}
        \scriptsize
        \begin{ganttchart}[
                canvas/.append style={draw=black!30},
                hgrid={draw=black!8},
                vgrid={draw=black!10},
                x unit=0.58cm,
                y unit chart=0.52cm,
                y unit title=0.5cm,
                title label font=\scriptsize\bfseries,
                title/.append style={fill=black!4, draw=black!20},
                bar label font=\scriptsize,
                milestone label font=\scriptsize,
                bar height=0.5,
                milestone/.append style={fill=ganttHito, draw=white, line width=0.4pt},
                milestone left shift=0.75,
                milestone right shift=0.25,
                milestone height=0.56,
            ]{1}{24}
            \gantttitle{\textcolor{ganttF1}{Fase~1 (M1--M6)}}{6}
            \gantttitle{\textcolor{ganttF2}{Fase~2 (M7--M12)}}{6}
            \gantttitle{\textcolor{ganttF3}{Fase~3 (M13--M24)}}{12} \\
            \gantttitlelist{1,2,3,4,5,6,7,8,9,10,11,12,13,14,15,16,17,18,19,20,21,22,23,24}{1} \\
            \ganttbar[bar/.append style={fill=ganttF1, draw=ganttF1!70!black}]{WP1.1 Inventario de identidades y modelo RBAC [H1.1]}{1}{3} \ganttmilestone{}{3} \\
            \ganttbar[bar/.append style={fill=ganttF1, draw=ganttF1!70!black}]{WP1.2 Activación progresiva de MFA [H1.4]}{3}{6} \ganttmilestone{}{6} \\
            \ganttbar[bar/.append style={fill=ganttF1, draw=ganttF1!70!black}]{WP1.3 Consolidación de RBAC y segregación}{4}{6} \\
            \ganttbar[bar/.append style={fill=ganttF1, draw=ganttF1!70!black}]{WP1.4 Registro automático de decisiones [H1.3]}{2}{5} \ganttmilestone{}{5} \\
            \ganttbar[bar/.append style={fill=ganttF1, draw=ganttF1!70!black}]{WP1.5 Fuentes autoritativas de datos [H1.2]}{1}{4} \ganttmilestone{}{4} \\
            \ganttbar[bar/.append style={fill=ganttF1, draw=ganttF1!70!black}]{WP1.6 Segmentación y protección perimetral}{4}{6} \\
            \ganttbar[bar/.append style={fill=ganttF1, draw=ganttF1!70!black}]{WP1.7 Verificación y cierre de Fase 1 [H1.5]}{6}{6} \ganttmilestone{}{6} \\
            \ganttbar[bar/.append style={fill=ganttF2, draw=ganttF2!70!black}]{WP2.1 Identidad para dispositivos ITS [H2.1]}{7}{9} \ganttmilestone{}{9} \\
            \ganttbar[bar/.append style={fill=ganttF2, draw=ganttF2!70!black}]{WP2.2 Autenticación de mensajes GPS/AVL}{9}{11} \\
            \ganttbar[bar/.append style={fill=ganttF2, draw=ganttF2!70!black}]{WP2.3 Plausibilidad y cuarentena de tramas [H2.2]}{9}{11} \ganttmilestone{}{11} \\
            \ganttbar[bar/.append style={fill=ganttF2, draw=ganttF2!70!black}]{WP2.4 Integración segura de bus y fuentes}{7}{11} \\
            \ganttbar[bar/.append style={fill=ganttF2, draw=ganttF2!70!black}]{WP2.5 WORM, sellado de tiempo y hashes [H2.3]}{7}{11} \ganttmilestone{}{11} \\
            \ganttbar[bar/.append style={fill=ganttF2, draw=ganttF2!70!black}]{WP2.6 Correlación SIEM y monitoreo}{9}{12} \\
            \ganttbar[bar/.append style={fill=ganttF2, draw=ganttF2!70!black}]{WP2.7 Verificación y cierre de Fase 2 [H2.4]}{12}{12} \ganttmilestone{}{12} \\
            \ganttbar[bar/.append style={fill=ganttF3, draw=ganttF3!70!black}]{WP3.1 Diseño de arquitectura híbrida [H3.1]}{13}{15} \ganttmilestone{}{15} \\
            \ganttbar[bar/.append style={fill=ganttF3, draw=ganttF3!70!black}]{WP3.2 Replicación de servicios y datos}{16}{20} \\
            \ganttbar[bar/.append style={fill=ganttF3, draw=ganttF3!70!black}]{WP3.3 Balanceo y failover automático}{18}{21} \\
            \ganttbar[bar/.append style={fill=ganttF3, draw=ganttF3!70!black}]{WP3.4 Consolidación del sitio alterno [H3.3]}{16}{21} \ganttmilestone{}{21} \\
            \ganttbar[bar/.append style={fill=ganttF3, draw=ganttF3!70!black}]{WP3.5 Formulación y aprobación de DRP/BCP [H3.2]}{13}{18} \ganttmilestone{}{18} \\
            \ganttbar[bar/.append style={fill=ganttF3, draw=ganttF3!70!black}]{WP3.6 Pruebas de recuperación y failover [H3.4]}{21}{23} \ganttmilestone{}{23} \\
            \ganttmilestone[milestone left shift=0.5, milestone right shift=0]{WP3.7 Verificación y cierre de Fase 3 [H3.5]}{24}
        \end{ganttchart}

        \vspace{2pt}
        {\footnotesize
            \raisebox{0.15ex}{\textcolor{ganttF1}{\rule{0.8em}{0.8em}}}\,Fase~1\qquad
            \raisebox{0.15ex}{\textcolor{ganttF2}{\rule{0.8em}{0.8em}}}\,Fase~2\qquad
            \raisebox{0.15ex}{\textcolor{ganttF3}{\rule{0.8em}{0.8em}}}\,Fase~3\qquad
            \raisebox{0.15ex}{\rotatebox[origin=c]{45}{\textcolor{ganttHito}{\rule{0.55em}{0.55em}}}}\,Hito
            \par}

        \vspace{2pt}
        \begin{minipage}{0.95\linewidth}
            \footnotesize
            \textit{Nota.} Las barras representan la duración de cada paquete de trabajo; los rombos rojos indican los hitos
            (identificador entre corchetes, véase la Tabla~\ref{tab:hitos_cci}) ubicados al cierre del paquete vinculado.
            Elaboración propia a partir del cronograma de la Tabla~\ref{tab:cronograma_wp}.
        \end{minipage}
    \end{figure}
    \endgroup
\end{landscape}


La Tabla~\ref{tab:hitos_cci} lista los 14 hitos del plan. Cada fase cierra con una puerta de verificación (H1.5, H2.4 y H3.5)
que corresponde al paso de un estado de transición al siguiente. Esta condición es coherente con el criterio establecido para la
hoja de ruta: una fase no se considera concluida por la instalación de una tecnología, sino cuando sus entregables han sido
verificados y la capacidad puede usarse de manera repetible, monitoreada y evaluada.

\begingroup
\scriptsize
\setlength{\tabcolsep}{3pt}
\renewcommand{\arraystretch}{1.15}
\begin{longtable}{
        >{\raggedright\arraybackslash}p{0.07\textwidth}
        >{\raggedright\arraybackslash}p{0.23\textwidth}
        >{\centering\arraybackslash}p{0.06\textwidth}
        >{\centering\arraybackslash}p{0.11\textwidth}
        >{\raggedright\arraybackslash}p{0.32\textwidth}
        >{\raggedright\arraybackslash}p{0.12\textwidth}}
    \caption{Hitos de la hoja de ruta y criterios de verificación}
    \label{tab:hitos_cci}                                                                                                                                                                                                                                                                \\
    \toprule
    \textbf{ID} & \textbf{Hito}                                     & \textbf{Mes} & \textbf{Entregable} & \textbf{Criterio de verificación}                                                                                                                       & \textbf{Aceptación} \\
    \midrule
    \endfirsthead
    \toprule
    \textbf{ID} & \textbf{Hito}                                     & \textbf{Mes} & \textbf{Entregable} & \textbf{Criterio de verificación}                                                                                                                       & \textbf{Aceptación} \\
    \midrule
    \endhead
    \midrule
    \multicolumn{6}{r}{\footnotesize Continúa en la página siguiente}                                                                                                                                                                                                                    \\
    \endfoot
    \bottomrule
    \endlastfoot
    \multicolumn{6}{l}{\textbf{Fase~1 · Transición 1}}                                                                                                                                                                                                                                   \\
    \midrule
    H1.1        & Matriz RBAC aprobada                              & M3           & E1.1                & Matriz de roles y permisos aprobada por la Dirección TIC y la Subgerencia de Operaciones, con mínimo privilegio y segregación de funciones definidos.   & Dirección TIC       \\
    H1.2        & Fuentes autoritativas de datos de flota definidas & M4           & E3.2                & Catálogo de fuentes autoritativas formalizado, con un responsable de dato asignado por cada fuente.                                                     & Datos               \\
    H1.3        & Registro automático de decisiones en operación    & M5           & E3.1                & Las decisiones del Sistema de Despacho se registran de forma automática en la bitácora, sin intervención manual.                                        & TIC                 \\
    H1.4        & MFA obligatorio en todas las consolas del CCI     & M6           & E1.2                & Ninguna consola del CCI permite el acceso sin segundo factor de autenticación.                                                                          & Seguridad           \\
    H1.5        & Cierre de Fase 1 – Transición 1                   & M6           & E1                  & Entregables E1, E3.1, E3.2 y E4.1 verificados contra sus criterios de aceptación; acta de cierre de fase aprobada.                                      & Dirección TIC       \\
    \addlinespace[3pt]
    \multicolumn{6}{l}{\textbf{Fase~2 · Transición 2}}                                                                                                                                                                                                                                   \\
    \midrule
    H2.1        & Identidad de dispositivos ITS desplegada          & M9           & E2.1                & Los dispositivos ITS dentro del alcance de la fase cuentan con identidad y certificado válidos, con ciclo de vida de certificados gestionado.           & Seguridad           \\
    H2.2        & Telemetría autenticada y validada en la pasarela  & M11          & E2.2                & La pasarela de ingreso rechaza o envía a cuarentena las tramas no autenticadas o con posiciones implausibles antes de que lleguen al flujo operacional. & Seguridad           \\
    H2.3        & Bitácora inmutable operativa                      & M11          & E3.3                & Registros protegidos con WORM, sellado de tiempo y encadenamiento de hashes; los intentos de modificación generan alerta.                               & Seguridad           \\
    H2.4        & Cierre de Fase 2 – Transición 2                   & M12          & E2                  & Entregables E2 y E3 verificados; el SIEM correlaciona eventos de telemetría y bitácora; acta de cierre de fase aprobada.                                & Dirección TIC       \\
    \addlinespace[3pt]
    \multicolumn{6}{l}{\textbf{Fase~3 · TO-BE}}                                                                                                                                                                                                                                          \\
    \midrule
    H3.1        & Arquitectura híbrida redundante aprobada          & M15          & E4.2                & Diseño CCI–GCP aprobado, con componentes, flujos de replicación y criterios de conmutación definidos.                                                   & Dirección TIC       \\
    H3.2        & DRP y BCP aprobados                               & M18          & E4.4                & DRP y BCP formalizados y aprobados, con objetivos de recuperación (RTO/RPO) definidos.                                                                  & Dirección TIC       \\
    H3.3        & Replicación y sitio alterno operativos            & M21          & E4.3                & Servicios y datos críticos replicados; sitio alterno consolidado y disponible para conmutación.                                                         & TIC                 \\
    H3.4        & Failover automático validado en pruebas           & M23          & E4.5                & Pruebas de recuperación y conmutación ejecutadas, documentadas y evaluadas contra los objetivos definidos en el DRP.                                    & TIC                 \\
    H3.5        & Cierre de Fase 3 – Estado TO-BE                   & M24          & E4                  & Entregable E4 verificado; puntos únicos de falla eliminados; estado TO-BE declarado con evidencia.                                                      & Dirección TIC       \\
\end{longtable}
\endgroup

{\footnotesize\textit{Nota.} Cada hito se ubica al cierre del paquete de trabajo vinculado (Tabla~\ref{tab:cronograma_wp}).
    Los hitos H1.5, H2.4 y H3.5 son las puertas de cierre de fase y corresponden a los estados Transición~1, Transición~2 y TO-BE.
    Los criterios de verificación son propuestas de esta hoja de ruta.\par}



\subsubsection{Entregables}

Los paquetes de trabajo producen los cuatro entregables principales definidos para la arquitectura de transición: \textit{Control de
    acceso seguro} (E1), \textit{Telemetría GPS confiable} (E2), \textit{Trazabilidad de los datos} (E3) y \textit{Disponibilidad y
    recuperación} (E4). Cada uno se descompone en productos de trabajo con un criterio de aceptación y un responsable, como muestra la
Tabla~\ref{tab:entregables_cci}. Algunos entregables maduran de forma incremental: los registros automáticos de decisiones y las
fuentes autoritativas (E3.1 y E3.2) y la segmentación perimetral (E4.1) se entregan en la Fase~1, aunque la trazabilidad completa
solo se consolida en la Fase~2 y la disponibilidad y recuperación en la Fase~3.

\begingroup
\scriptsize
\setlength{\tabcolsep}{3pt}
\renewcommand{\arraystretch}{1.15}
\begin{longtable}{
        >{\raggedright\arraybackslash}p{0.07\textwidth}
        >{\raggedright\arraybackslash}p{0.25\textwidth}
        >{\raggedright\arraybackslash}p{0.07\textwidth}
        >{\centering\arraybackslash}p{0.05\textwidth}
        >{\raggedright\arraybackslash}p{0.35\textwidth}
        >{\raggedright\arraybackslash}p{0.12\textwidth}}
    \caption{Entregables principales y productos de trabajo de la hoja de ruta}
    \label{tab:entregables_cci}                                                                                                                                                                                                                                                                 \\
    \toprule
    \textbf{ID} & \textbf{Entregable}                                                               & \textbf{Fase} & \textbf{Mes} & \textbf{Criterio de aceptación}                                                                                                     & \textbf{Responsable} \\
    \midrule
    \endfirsthead
    \toprule
    \textbf{ID} & \textbf{Entregable}                                                               & \textbf{Fase} & \textbf{Mes} & \textbf{Criterio de aceptación}                                                                                                     & \textbf{Responsable} \\
    \midrule
    \endhead
    \midrule
    \multicolumn{6}{r}{\footnotesize Continúa en la página siguiente}                                                                                                                                                                                                                           \\
    \endfoot
    \bottomrule
    \endlastfoot
    \textbf{E1} & \textbf{Control de acceso seguro}                                                 & \textbf{F1}   & \textbf{M6}  & \textbf{MFA obligatorio activo en las consolas del CCI y matriz RBAC aprobada y aplicada, con segregación de funciones verificada.} & \textbf{Seguridad}   \\
    E1.1        & \hspace{0.6em}Matriz RBAC consolidada                                             & F1            & M6           & Roles, permisos y segregación de funciones documentados, aprobados y aplicados en las consolas.                                     & Seguridad            \\
    E1.2        & \hspace{0.6em}MFA obligatorio en las consolas del CCI                             & F1            & M6           & Segundo factor exigido en todas las consolas; accesos sin MFA bloqueados.                                                           & Seguridad            \\
    \addlinespace[3pt]
    \textbf{E2} & \textbf{Telemetría GPS confiable}                                                 & \textbf{F2}   & \textbf{M12} & \textbf{Solo la telemetría autenticada y validada llega al Sistema de Despacho, integrada con fuentes autoritativas.}               & \textbf{Seguridad}   \\
    E2.1        & \hspace{0.6em}Identidad y certificados de dispositivos ITS                        & F2            & M9           & Cada dispositivo ITS en alcance se identifica con credencial propia y vigente.                                                      & Seguridad            \\
    E2.2        & \hspace{0.6em}Pasarela de ingreso con autenticación y validación de plausibilidad & F2            & M11          & Tramas inválidas rechazadas o en cuarentena; mensajes válidos con autenticidad e integridad verificadas.                            & Seguridad            \\
    E2.3        & \hspace{0.6em}Integración segura del bus de eventos y fuentes geoespaciales       & F2            & M11          & Flujos entre telemetría, bus de eventos y fuentes autoritativas integrados con canales y accesos controlados.                       & TIC                  \\
    \addlinespace[3pt]
    \textbf{E3} & \textbf{Trazabilidad de los datos}                                                & \textbf{F2}   & \textbf{M12} & \textbf{Cada decisión operacional es rastreable hasta su fuente de datos y su autor, con evidencia protegida contra alteración.}    & \textbf{Seguridad}   \\
    E3.1        & \hspace{0.6em}Registro automático de decisiones                                   & F1            & M5           & Decisiones registradas de forma automática en la bitácora de operación.                                                             & TIC                  \\
    E3.2        & \hspace{0.6em}Catálogo de fuentes autoritativas de datos de flota                 & F1            & M4           & Fuentes identificadas y formalizadas, con propietario de dato por fuente.                                                           & Datos                \\
    E3.3        & \hspace{0.6em}Bitácora inmutable (WORM, sellado de tiempo y hashes)               & F2            & M11          & Registros no modificables, con sello de tiempo y encadenamiento verificable; alertas ante intentos de alteración.                   & Seguridad            \\
    E3.4        & \hspace{0.6em}Correlación SIEM y monitoreo centralizado                           & F2            & M12          & Eventos de telemetría, accesos y bitácora correlacionados y monitoreados de forma centralizada.                                     & Seguridad            \\
    \addlinespace[3pt]
    \textbf{E4} & \textbf{Disponibilidad y recuperación}                                            & \textbf{F3}   & \textbf{M24} & \textbf{Arquitectura híbrida redundante en operación y DRP probado; sin puntos únicos de falla.}                                    & \textbf{TIC}         \\
    E4.1        & \hspace{0.6em}Segmentación y protección perimetral inicial                        & F1            & M6           & Dominios segmentados y reglas de flujo entre ellos aplicadas en la frontera.                                                        & Seguridad            \\
    E4.2        & \hspace{0.6em}Arquitectura híbrida redundante con replicación                     & F3            & M20          & Servicios y datos críticos replicados entre infraestructura del CCI y GCP.                                                          & TIC                  \\
    E4.3        & \hspace{0.6em}Conmutación automática y sitio alterno                              & F3            & M21          & Conmutación automática funcional hacia un sitio alterno consolidado.                                                                & TIC                  \\
    E4.4        & \hspace{0.6em}DRP y BCP aprobados                                                 & F3            & M18          & Planes formalizados y aprobados, con roles, procedimientos y objetivos de recuperación.                                             & Dirección TIC        \\
    E4.5        & \hspace{0.6em}Informe de pruebas de recuperación y failover                       & F3            & M23          & Pruebas ejecutadas y documentadas, con resultados evaluados contra los objetivos del DRP.                                           & TIC                  \\
\end{longtable}
\endgroup

{\footnotesize\textit{Nota.} El mes de cada entregable principal corresponde al hito de cierre de la fase en que se consolida;
    el de cada producto de trabajo, al cierre del paquete que lo produce. La trazabilidad de los datos (E3) se consolida en la
    Fase~2, pero sus primeros productos (E3.1 y E3.2) se entregan en la Fase~1, y la disponibilidad y recuperación (E4)
    inicia con la segmentación perimetral (E4.1) en la Fase~1. Elaboración propia.\par}



\subsubsection{Asignación de recursos}

La ejecución involucra siete grupos de recursos humanos: la Dirección TIC, el Equipo de Seguridad, el Equipo TIC, la Subgerencia de
Operaciones, los responsables de los datos operacionales críticos, los concesionarios de flota y la función de auditoría o control
interno. El documento identifica un único «Equipo de Seguridad y TIC»; para poder asignar responsabilidades y estimar esfuerzo, esta
hoja de ruta lo desagrega en dos roles, Seguridad y TIC. La Tabla~\ref{tab:raci_cci} asigna para cada paquete de trabajo un único
aprobador (A) y define quién ejecuta, quién es consultado y quién debe ser informado.

\begin{table}[H]
    \centering
    \caption{Matriz RACI por paquete de trabajo}
    \label{tab:raci_cci}
    \scriptsize
    \setlength{\tabcolsep}{3pt}
    \renewcommand{\arraystretch}{1.15}

    \begin{tabularx}{\textwidth}{
            >{\raggedright\arraybackslash}p{1.0cm}
            >{\raggedright\arraybackslash}X
                                              *{7}{>{\centering\arraybackslash}p{0.95cm}}}
        \toprule
        \textbf{ID} & \textbf{Paquete de trabajo}                                   & \textbf{DTIC} & \textbf{SEG} & \textbf{TIC} & \textbf{OPS} & \textbf{DAT} & \textbf{CON} & \textbf{AUD} \\
        \midrule
        \multicolumn{9}{l}{\textbf{Fase~1}}                                                                                                                                                   \\
        \midrule
        WP1.1       & Inventario de identidades y diseño del modelo RBAC            & I             & A            & R            & C            &              &              & I            \\
        WP1.2       & Activación progresiva de MFA en las consolas del CCI          & I             & A            & R            & C            &              &              & I            \\
        WP1.3       & Consolidación de RBAC y segregación de funciones              & I             & A            & R            & C            &              &              & C            \\
        WP1.4       & Registro automático de decisiones                             & I             & C            & A/R          & C            & C            &              & I            \\
        WP1.5       & Definición de fuentes autoritativas de datos de flota         & I             & C            & R            & C            & A/R          & C            & I            \\
        WP1.6       & Segmentación y protección perimetral inicial                  & I             & A            & R            & I            &              &              &              \\
        WP1.7       & Verificación de entregables y cierre de la Fase 1             & A             & R            & C            & C            & C            &              & R            \\
        \addlinespace[3pt]
        \multicolumn{9}{l}{\textbf{Fase~2}}                                                                                                                                                   \\
        \midrule
        WP2.1       & Identidad para dispositivos ITS                               & I             & A            & R            & I            &              & R            &              \\
        WP2.2       & Autenticación de mensajes GPS/AVL en la pasarela de ingreso   & I             & A            & R            &              & C            & C            &              \\
        WP2.3       & Validación de plausibilidad y cuarentena de tramas inválidas  &               & A/R          & R            & C            & C            & C            &              \\
        WP2.4       & Integración segura del bus de eventos y fuentes geoespaciales &               & C            & A/R          & I            & R            & C            &              \\
        WP2.5       & WORM, sellado de tiempo y encadenamiento de hashes            &               & A            & R            &              & C            &              & C            \\
        WP2.6       & Correlación en SIEM y monitoreo centralizado                  & I             & A/R          & R            & C            &              &              & I            \\
        WP2.7       & Verificación de entregables y cierre de la Fase 2             & A             & R            & C            & C            & C            &              & R            \\
        \addlinespace[3pt]
        \multicolumn{9}{l}{\textbf{Fase~3}}                                                                                                                                                   \\
        \midrule
        WP3.1       & Diseño de la arquitectura híbrida redundante                  & C             & R            & A            & C            & I            &              &              \\
        WP3.2       & Replicación de servicios y datos CCI–GCP                      & I             & C            & A/R          &              & C            &              &              \\
        WP3.3       & Balanceo y conmutación automática (failover)                  &               & C            & A/R          & I            &              &              &              \\
        WP3.4       & Consolidación del sitio alterno                               & I             & C            & A/R          & C            &              &              &              \\
        WP3.5       & Formulación y aprobación de DRP y BCP                         & A             & R            & R            & R            & C            &              & C            \\
        WP3.6       & Pruebas de recuperación y failover                            & I             & R            & A            & R            & C            &              & C            \\
        WP3.7       & Verificación de entregables y cierre de la Fase 3 (TO-BE)     & A             & R            & C            & C            & C            &              & R            \\
        \bottomrule
    \end{tabularx}

    \begin{minipage}{0.96\textwidth}
        \footnotesize
        \textit{Nota.} R: responsable de ejecutar; A: aprobador que rinde cuentas (uno por paquete); C: consultado; I: informado;
        A/R: ambos roles. DTIC: Dirección TIC; SEG: Equipo de Seguridad; TIC: Equipo TIC (infraestructura y plataformas);
        OPS: Subgerencia de Operaciones; DAT: responsables de datos operacionales críticos; CON: concesionarios de flota;
        AUD: auditoría o control interno. Elaboración propia.
    \end{minipage}
\end{table}


En la asignación, el Equipo de Seguridad rinde cuentas de 9 de los 21 paquetes, principalmente los de identidad, telemetría y
trazabilidad; el Equipo TIC, de 7, principalmente los de integración, redundancia y recuperación; y la Dirección TIC, de 4,
que corresponden a las verificaciones de cierre de fase y a la aprobación de los planes DRP y BCP. Los responsables de datos operacionales críticos
lideran la definición de las fuentes autoritativas (WP1.5), y los concesionarios participan de forma directa en la identidad de los
dispositivos ITS (WP2.1), dado que intervienen los equipos instalados a bordo de la flota.

Como soporte tecnológico, la Fase~1 se apoya en Active Directory y ADFS, la plataforma IAM y el servicio MFA, el Sistema de Despacho y su
bitácora, las plataformas geoespaciales y de gestión de flota (Oracle Spatial/Locator, ArcGIS y SIRCI/FMS) y los firewalls Fortinet.
La Fase~2 requiere servicios criptográficos y gestor de claves, los equipos ITS a bordo, la pasarela de ingreso con TLS, el bus de eventos
y las APIs, el almacenamiento WORM con servicio de hora y el SIEM. La Fase~3 se apoya en la infraestructura tecnológica del CCI, Google
Cloud Platform, los balanceadores de carga y firewalls, y las plataformas de continuidad y respaldo.

La Tabla~\ref{tab:esfuerzo_cci} resume la estimación preliminar de esfuerzo. El esfuerzo total estimado del programa es de 84{,}6 persona-mes: el
Equipo TIC aporta 43{,}6\,\% del esfuerzo y el Equipo de Seguridad 28{,}7\,\%. La carga del Equipo TIC se concentra en la Fase~3
    (16{,}7 persona-mes), coherente con el carácter de infraestructura de la redundancia híbrida y la recuperación, mientras que la
participación de los concesionarios se concentra en la Fase~2, cuando se incorporan identidades a los dispositivos ITS.

\begin{table}[H]
    \centering
    \caption{Estimación preliminar de esfuerzo por fase y rol (persona-mes)}
    \label{tab:esfuerzo_cci}
    \small
    \setlength{\tabcolsep}{4pt}
    \renewcommand{\arraystretch}{1.2}

    \begin{tabularx}{\textwidth}{
            >{\raggedright\arraybackslash}p{2.0cm}
                                              *{7}{>{\centering\arraybackslash}X}
            >{\centering\arraybackslash}p{1.4cm}
            >{\centering\arraybackslash}p{1.5cm}}
        \toprule
        \textbf{Fase}     & \textbf{DTIC}  & \textbf{SEG}    & \textbf{TIC}    & \textbf{OPS}   & \textbf{DAT}   & \textbf{CON}   & \textbf{AUD}   & \textbf{Total}  & \textbf{FTE medio} \\
        \midrule
        Fase~1 (M1--M6)   & {,}8           & 8{,}9           & 9{,}0           & 3{,}6          & 2{,}0          & {,}4           & {,}7           & 25{,}4          & 4{,}2              \\
        Fase~2 (M7--M12)  & {,}2           & 8{,}9           & 11{,}2          & 1{,}5          & 2{,}6          & 3{,}8          & {,}8           & 29{,}0          & 4{,}8              \\
        Fase~3 (M13--M24) & 1{,}1          & 6{,}5           & 16{,}7          & 3{,}5          & 1{,}3          & {,}0           & 1{,}1          & 30{,}2          & 2{,}5              \\
        \midrule
        \textbf{Total}    & \textbf{2{,}1} & \textbf{24{,}3} & \textbf{36{,}9} & \textbf{8{,}6} & \textbf{5{,}9} & \textbf{4{,}2} & \textbf{2{,}6} & \textbf{84{,}6} & \textbf{3{,}5}     \\
        \bottomrule
    \end{tabularx}

    \begin{minipage}{0.96\textwidth}
        \footnotesize
        \textit{Nota.} Estimación preliminar de planificación, sujeta a validación por la Dirección TIC; no incluye costos de
        licencias, infraestructura ni contratación externa. El FTE medio resulta de dividir el esfuerzo de la fase entre su
        duración en meses. Los roles se identifican con las mismas siglas de la Tabla~\ref{tab:raci_cci}. Elaboración propia.
    \end{minipage}
\end{table}



\subsubsection{Dependencias, supuestos y condiciones de cierre}

El plan presenta dependencias entre fases que deben gestionarse de forma explícita. La definición de fuentes autoritativas (WP1.5)
condiciona la integración del bus de eventos y de las fuentes geoespaciales (WP2.4), y el registro automático de decisiones (WP1.4)
es la base sobre la cual se construye la bitácora inmutable (WP2.5). La correlación en SIEM (WP2.6) requiere telemetría autenticada
(WP2.2) y bitácora protegida (WP2.5), y el cierre de la Fase~2 (WP2.7) habilita el inicio del diseño de la arquitectura redundante
(WP3.1) y de los planes de continuidad (WP3.5). Un retraso en cualquiera de estos paquetes se propaga a las fases siguientes.

Los principales supuestos del plan son los siguientes. Primero, el documento no fija una fecha calendario de inicio, por lo que
el cronograma se expresa en meses relativos. Segundo, las duraciones y el esfuerzo son estimaciones preliminares y deben validarse
antes de su adopción. Tercero, el alcance de dispositivos ITS cubierto en cada fase debe acordarse con los concesionarios de flota. Cuarto, la
activación de MFA y la segmentación perimetral se realizan de forma progresiva para no comprometer la continuidad de la operación
del CCI.

Finalmente, la hoja de ruta no se limita a un cronograma tecnológico: cada paquete de trabajo se asocia con un entregable, un hito
verificable y un responsable, de modo que el avance pueda evaluarse contra los resultados esperados de la arquitectura objetivo.